\documentclass[11pt,a4paper]{article}

\usepackage[margin=2.3cm]{geometry}
% Overleaf-safe fonts (compile with the DEFAULT pdfLaTeX; no system fonts needed).
% XCharter ~ Charter (body); TeX Gyre Adventor ~ Avenir-style geometric sans
% (headings); mono for code. All ship with TeX Live / Overleaf.
\usepackage[T1]{fontenc}
\usepackage{XCharter}
\usepackage{tgadventor}
\usepackage[scaled=0.85]{beramono}
\newcommand{\headfont}{\sffamily}
\usepackage{amsmath}
\usepackage{enumitem}
\usepackage{xcolor}
\usepackage{titlesec}
\usepackage{parskip}
\usepackage{booktabs}
\usepackage{array}
\usepackage[colorlinks=true,linkcolor=ink,urlcolor=linkc]{hyperref}

\definecolor{ink}{HTML}{1A1A1A}
\definecolor{mid}{HTML}{444444}
\definecolor{muted}{HTML}{666666}
\definecolor{rule}{HTML}{CCCCCC}
\definecolor{linkc}{HTML}{2B5FA8}
\color{ink}

\titleformat{\section}{\headfont\large\bfseries\color{ink}}{\thesection}{0.8em}{}
\titlespacing{\section}{0pt}{0.8em}{0.25em}
\titleformat{\subsection}{\headfont\normalsize\bfseries\color{mid}}{\thesubsection}{0.6em}{}
\titlespacing{\subsection}{0pt}{0.5em}{0.2em}

\newcommand{\thinrule}{\noindent\textcolor{rule}{\rule{\linewidth}{0.6pt}}}
\setlist[itemize]{leftmargin=1.3em, itemsep=1pt, topsep=2pt, label=\textcolor{muted}{\textbullet}}
\setlength{\parskip}{0.35em}
\setlength{\emergencystretch}{5em}

\begin{document}

{\headfont\Large\bfseries How to read our results}
\hfill{\color{muted}\small Companion to \texttt{proposal.pdf} \,\textperiodcentered\, ACVSS Hackathon}

\vspace{0.25em}
\thinrule
\vspace{0.2em}

{\headfont\normalsize\bfseries\color{ink} Terms}\\[0.1em]
Read these first. Three are used more narrowly here than in ordinary usage, and every
confusing number below traces back to one of them.

\begin{itemize}\small
  \item \textbf{A lesion is} an area of abnormal tissue, so for us a tumor: one
        three-dimensional mass a radiologist would put in a report. \textbf{In our metrics a
        lesion is} something much narrower --- one connected blob of tumor mask within a
        \emph{single 2D slice}. A scan is a stack of about 155 cross-sections, so one tumor
        crossing twenty of them counts as twenty lesions, a thin or broken enhancing rim
        fragments into several more inside each slice, and a five-pixel speck counts the same
        as a whole tumor.
  \item \textbf{Erasure is} the failure we set out to fix.
        \textbf{The erasure rate is} the number of true lesions the
        model failed to recover, divided by the number of true lesions, where a lesion counts
        as recovered if the prediction covers at least 10\% of it. Lower is safer.
  \item \textbf{Hallucination is} the opposite failure: the reconstruction invents tumor where
        there is none. \textbf{The hallucination rate is} the number of predicted blobs that
        are fabricated, divided by the number of predicted blobs, where a blob counts as
        fabricated if less than 10\% of it sits inside real tumor. Lower is safer, and it tends
        to move against erasure.
  \item \textbf{The floor is} what the tumor detector misses when handed the untouched original
        scan, with no degradation or reconstruction involved. It is large, so it is the
        denominator every erasure rate must be read against.
  \item \textbf{The baseline is} the distortion-optimal model: the standard approach, trained
        only to minimise average pixel error. \textbf{Ours is} the tumor-aware model: the same
        network and the same data, with errors inside the tumor weighted far more heavily.
\end{itemize}

\vspace{0.3em}
\thinrule

\section*{The result}

Sharpening a cheap blurry scan with AI can quietly delete a small tumor, because the model is
graded on average pixel error and a small tumor is a tiny fraction of the pixels. We changed
what the model is punished for, weighting errors inside the tumor far more heavily, and then
measured how much tumor survives. On 70 held-out BraTS patients it loses far less tumor than
the standard model, at the \emph{same} segmentation quality.

\begin{center}\small
\renewcommand{\arraystretch}{1.15}
\begin{tabular}{@{}>{\raggedright\arraybackslash}p{6.4cm}cc@{}}
\toprule
\textbf{What the tumor detector was given} & \textbf{Tumor missed} & \textbf{Caused by sharpening} \\
\midrule
The original scan, untouched           & 45.5\% & --- \\
Standard sharpening (baseline)         & 54.6\% & $+9.1$ points \\
\textbf{Tumor-aware sharpening (ours)} & \textbf{48.1\%} & \textbf{$+2.6$ points} \\
\bottomrule
\end{tabular}
\end{center}

\textbf{Read the right-hand column, not the middle one.} The detector already misses 45.5\% on
a perfect scan, so most of the middle column is nothing to do with sharpening. Sharpening adds
9.1 points on top of that floor and \textbf{our objective removes 71\% of that damage}. 66 of
70 patients improve, 1 is worse, $p < 0.0001$. Quality is genuinely matched, not approximately:
Dice 0.660 against 0.675, and brain-masked PSNR within 0.23~dB.

\subsection*{Is a rate near 50\% bad? Split it by lesion size}

\begin{center}\small
\renewcommand{\arraystretch}{1.15}
\begin{tabular}{@{}lcccc@{}}
\toprule
\textbf{Lesion size} & \textbf{How many} & \textbf{Floor} & \textbf{Baseline adds} & \textbf{Ours adds} \\
\midrule
Small, $<50$ px    & 6{,}762 (71\%) & 65.1\% & $+13.5$ & $+4.7$ \\
Medium, 50--200 px & 1{,}005        & 10.0\% & $+6.4$  & $+3.5$ \\
Large, $>200$ px   & 1{,}723        & 0.5\%  & $+0.7$  & $+0.4$ \\
\bottomrule
\end{tabular}
\end{center}

Read the bottom row first: a lesion above 200~px is missed 0.5\% of the time on a perfect scan
and 0.9\% after our reconstruction, so substantial tumors are not vanishing. The overall rate
is high because 71\% of components are under 50~px and two thirds of those are missed even on
the original image. Small does not mean unimportant --- a 30~px enhancing focus can be an early
recurrence, and this is where the remaining work is --- but it is also partly an artefact of
counting a fragmenting rim as many equal lesions.

\subsection*{Three things that make it credible, and one that limits it}

\begin{itemize}
  \item \textbf{It appears only where the theory says it must.} The benefit is there for the
        small enhancing tumor and \emph{vanishes} for the whole tumor region, nine times
        larger. That is the prediction: the objective is only misaligned when the structure is
        too small to matter to a pixel-error score.
  \item \textbf{The patient is the unit of analysis.} One tumor spans many slices and appears
        as a lesion in each, so those outcomes are correlated; pooling them fakes an interval
        about 1.5 times too tight. Resampling patients instead gives a mean paired difference
        of 6.5 points, 95\% CI $[5.2, 8.2]$.
  \item \textbf{The weakest link is the detector, not the sharpening.} It misses 45.5\% of
        lesions on an untouched scan. A stronger downstream segmenter would buy more than a
        better loss function.
  \item \textbf{The limit: it is not free.} Protecting tumors makes the model readier to report
        one that is not there, and hallucination rises from 0.266 to 0.387. A second loss term
        recovers most of that cost for about half the benefit, so the tradeoff is a dial rather
        than a fixed price.
\end{itemize}

\vspace{0.2em}
{\small\itshape Provenance: enhancing-tumor checkpoint trained on 17{,}233 axial slices from
468 BraTS patients (Medical Segmentation Decathlon brain task), split by patient, measured on
the 70-patient validation split. Raw outputs in \texttt{results/}; checkpoints at
\texttt{douyeszn/tumor-aware-sr} on Hugging Face. The proposal's results table is from an
earlier, smaller run, so its digits differ.}

\end{document}
